\batchmode
\documentclass[12pt]{article}
\RequirePackage{ifthen}


\usepackage{amsmath}
\usepackage{mathrsfs}
\usepackage{eucal}
\usepackage{amssymb}
\usepackage[french]{babel}
\usepackage[latin1]{inputenc}
\usepackage[T1]{fontenc}
\usepackage[all]{xy}
\usepackage{html}

%
\providecommand{\perm}[3]{{\rm signe} \left( \begin{array}{ccc} 1 & 2 & 3 \\#1 & #2 & #3 \end{array} \right) \cdot a_{1,#1} a_{2,#2} a_{3,#3} }%
\providecommand{\perma}[3]{ \cdot a_{1,#1} a_{2,#2} a_{3,#3} }%
\providecommand{\semblableun}[1]{\overset{{\tiny\textrm{$\begin{array}{c} #1 \end{array}$}}}{\sim}}%
\providecommand{\semblabledeux}[2]{\overset{{\tiny\textrm{$\begin{array}{c} #1 \\#2 \end{array}$}}}{\sim}}%
\providecommand{\semblabletrois}[3]{\overset{\begin{array}{c} #1 \\#2 \\#3 \end{array}}{\sim}} 




\usepackage[dvips]{color}


\pagecolor[gray]{.7}

\usepackage[latin1]{inputenc}



\makeatletter

\makeatletter
\count@=\the\catcode`\_ \catcode`\_=8 
\newenvironment{tex2html_wrap}{}{}%
\catcode`\<=12\catcode`\_=\count@
\newcommand{\providedcommand}[1]{\expandafter\providecommand\csname #1\endcsname}%
\newcommand{\renewedcommand}[1]{\expandafter\providecommand\csname #1\endcsname{}%
  \expandafter\renewcommand\csname #1\endcsname}%
\newcommand{\newedenvironment}[1]{\newenvironment{#1}{}{}\renewenvironment{#1}}%
\let\newedcommand\renewedcommand
\let\renewedenvironment\newedenvironment
\makeatother
\let\mathon=$
\let\mathoff=$
\ifx\AtBeginDocument\undefined \newcommand{\AtBeginDocument}[1]{}\fi
\newbox\sizebox
\setlength{\hoffset}{0pt}\setlength{\voffset}{0pt}
\addtolength{\textheight}{\footskip}\setlength{\footskip}{0pt}
\addtolength{\textheight}{\topmargin}\setlength{\topmargin}{0pt}
\addtolength{\textheight}{\headheight}\setlength{\headheight}{0pt}
\addtolength{\textheight}{\headsep}\setlength{\headsep}{0pt}
\setlength{\textwidth}{349pt}
\newwrite\lthtmlwrite
\makeatletter
\let\realnormalsize=\normalsize
\global\topskip=2sp
\def\preveqno{}\let\real@float=\@float \let\realend@float=\end@float
\def\@float{\let\@savefreelist\@freelist\real@float}
\def\liih@math{\ifmmode$\else\bad@math\fi}
\def\end@float{\realend@float\global\let\@freelist\@savefreelist}
\let\real@dbflt=\@dbflt \let\end@dblfloat=\end@float
\let\@largefloatcheck=\relax
\let\if@boxedmulticols=\iftrue
\def\@dbflt{\let\@savefreelist\@freelist\real@dbflt}
\def\adjustnormalsize{\def\normalsize{\mathsurround=0pt \realnormalsize
 \parindent=0pt\abovedisplayskip=0pt\belowdisplayskip=0pt}%
 \def\phantompar{\csname par\endcsname}\normalsize}%
\def\lthtmltypeout#1{{\let\protect\string \immediate\write\lthtmlwrite{#1}}}%
\newcommand\lthtmlhboxmathA{\adjustnormalsize\setbox\sizebox=\hbox\bgroup\kern.05em }%
\newcommand\lthtmlhboxmathB{\adjustnormalsize\setbox\sizebox=\hbox to\hsize\bgroup\hfill }%
\newcommand\lthtmlvboxmathA{\adjustnormalsize\setbox\sizebox=\vbox\bgroup %
 \let\ifinner=\iffalse \let\)\liih@math }%
\newcommand\lthtmlboxmathZ{\@next\next\@currlist{}{\def\next{\voidb@x}}%
 \expandafter\box\next\egroup}%
\newcommand\lthtmlmathtype[1]{\gdef\lthtmlmathenv{#1}}%
\newcommand\lthtmllogmath{\dimen0\ht\sizebox \advance\dimen0\dp\sizebox
  \ifdim\dimen0>.95\vsize
   \lthtmltypeout{%
*** image for \lthtmlmathenv\space is too tall at \the\dimen0, reducing to .95 vsize ***}%
   \ht\sizebox.95\vsize \dp\sizebox\z@ \fi
  \lthtmltypeout{l2hSize %
:\lthtmlmathenv:\the\ht\sizebox::\the\dp\sizebox::\the\wd\sizebox.\preveqno}}%
\newcommand\lthtmlfigureA[1]{\let\@savefreelist\@freelist
       \lthtmlmathtype{#1}\lthtmlvboxmathA}%
\newcommand\lthtmlpictureA{\bgroup\catcode`\_=8 \lthtmlpictureB}%
\newcommand\lthtmlpictureB[1]{\lthtmlmathtype{#1}\egroup
       \let\@savefreelist\@freelist \lthtmlhboxmathB}%
\newcommand\lthtmlpictureZ[1]{\hfill\lthtmlfigureZ}%
\newcommand\lthtmlfigureZ{\lthtmlboxmathZ\lthtmllogmath\copy\sizebox
       \global\let\@freelist\@savefreelist}%
\newcommand\lthtmldisplayA{\bgroup\catcode`\_=8 \lthtmldisplayAi}%
\newcommand\lthtmldisplayAi[1]{\lthtmlmathtype{#1}\egroup\lthtmlvboxmathA}%
\newcommand\lthtmldisplayB[1]{\edef\preveqno{(\theequation)}%
  \lthtmldisplayA{#1}\let\@eqnnum\relax}%
\newcommand\lthtmldisplayZ{\lthtmlboxmathZ\lthtmllogmath\lthtmlsetmath}%
\newcommand\lthtmlinlinemathA{\bgroup\catcode`\_=8 \lthtmlinlinemathB}
\newcommand\lthtmlinlinemathB[1]{\lthtmlmathtype{#1}\egroup\lthtmlhboxmathA
  \vrule height1.5ex width0pt }%
\newcommand\lthtmlinlineA{\bgroup\catcode`\_=8 \lthtmlinlineB}%
\newcommand\lthtmlinlineB[1]{\lthtmlmathtype{#1}\egroup\lthtmlhboxmathA}%
\newcommand\lthtmlinlineZ{\egroup\expandafter\ifdim\dp\sizebox>0pt %
  \expandafter\centerinlinemath\fi\lthtmllogmath\lthtmlsetinline}
\newcommand\lthtmlinlinemathZ{\egroup\expandafter\ifdim\dp\sizebox>0pt %
  \expandafter\centerinlinemath\fi\lthtmllogmath\lthtmlsetmath}
\newcommand\lthtmlindisplaymathZ{\egroup %
  \centerinlinemath\lthtmllogmath\lthtmlsetmath}
\def\lthtmlsetinline{\hbox{\vrule width.1em \vtop{\vbox{%
  \kern.1em\copy\sizebox}\ifdim\dp\sizebox>0pt\kern.1em\else\kern.3pt\fi
  \ifdim\hsize>\wd\sizebox \hrule depth1pt\fi}}}
\def\lthtmlsetmath{\hbox{\vrule width.1em\kern-.05em\vtop{\vbox{%
  \kern.1em\kern0.8 pt\hbox{\hglue.17em\copy\sizebox\hglue0.8 pt}}\kern.3pt%
  \ifdim\dp\sizebox>0pt\kern.1em\fi \kern0.8 pt%
  \ifdim\hsize>\wd\sizebox \hrule depth1pt\fi}}}
\def\centerinlinemath{%
  \dimen1=\ifdim\ht\sizebox<\dp\sizebox \dp\sizebox\else\ht\sizebox\fi
  \advance\dimen1by.5pt \vrule width0pt height\dimen1 depth\dimen1 
 \dp\sizebox=\dimen1\ht\sizebox=\dimen1\relax}

\def\lthtmlcheckvsize{\ifdim\ht\sizebox<\vsize 
  \ifdim\wd\sizebox<\hsize\expandafter\hfill\fi \expandafter\vfill
  \else\expandafter\vss\fi}%
\providecommand{\selectlanguage}[1]{}%
\makeatletter \tracingstats = 1 
\providecommand{\Beta}{\textrm{B}}
\providecommand{\Mu}{\textrm{M}}
\providecommand{\Kappa}{\textrm{K}}
\providecommand{\Rho}{\textrm{R}}
\providecommand{\Epsilon}{\textrm{E}}
\providecommand{\Chi}{\textrm{X}}
\providecommand{\Iota}{\textrm{J}}
\providecommand{\omicron}{\textrm{o}}
\providecommand{\Zeta}{\textrm{Z}}
\providecommand{\Eta}{\textrm{H}}
\providecommand{\Nu}{\textrm{N}}
\providecommand{\Omicron}{\textrm{O}}
\providecommand{\Tau}{\textrm{T}}
\providecommand{\Alpha}{\textrm{A}}


\begin{document}
\pagestyle{empty}\thispagestyle{empty}\lthtmltypeout{}%
\lthtmltypeout{latex2htmlLength hsize=\the\hsize}\lthtmltypeout{}%
\lthtmltypeout{latex2htmlLength vsize=\the\vsize}\lthtmltypeout{}%
\lthtmltypeout{latex2htmlLength hoffset=\the\hoffset}\lthtmltypeout{}%
\lthtmltypeout{latex2htmlLength voffset=\the\voffset}\lthtmltypeout{}%
\lthtmltypeout{latex2htmlLength topmargin=\the\topmargin}\lthtmltypeout{}%
\lthtmltypeout{latex2htmlLength topskip=\the\topskip}\lthtmltypeout{}%
\lthtmltypeout{latex2htmlLength headheight=\the\headheight}\lthtmltypeout{}%
\lthtmltypeout{latex2htmlLength headsep=\the\headsep}\lthtmltypeout{}%
\lthtmltypeout{latex2htmlLength parskip=\the\parskip}\lthtmltypeout{}%
\lthtmltypeout{latex2htmlLength oddsidemargin=\the\oddsidemargin}\lthtmltypeout{}%
\makeatletter
\if@twoside\lthtmltypeout{latex2htmlLength evensidemargin=\the\evensidemargin}%
\else\lthtmltypeout{latex2htmlLength evensidemargin=\the\oddsidemargin}\fi%
\lthtmltypeout{}%
\makeatother
\setcounter{page}{1}
\onecolumn

% !!! IMAGES START HERE !!!

{\newpage\clearpage
\lthtmldisplayA{displaymath1276}%
\begin{displaymath}\begin{array}[t]{rcl}
\det A  & = & \sum\limits_{\sigma \in \mathcal{S}_3} \left( {\rm sgn}(\sigma) \prod\limits_{i=1}^{3} a_{i, \sigma(i)} \right) \\\\
        & = & {\rm signe} \left( \begin{array}{ccc} 1 & 2 & 3 \\1 & 2 & 3 \end{array} \right) \cdot a_{1,1} a_{2,2} a_{3,3}  + {\rm signe} \left( \begin{array}{ccc} 1 & 2 & 3 \\1 & 3 & 2 \end{array} \right) \cdot a_{1,1} a_{2,3} a_{3,2}   \\\\
        &   & \quad + {\rm signe} \left( \begin{array}{ccc} 1 & 2 & 3 \\2 & 1 & 3 \end{array} \right) \cdot a_{1,2} a_{2,1} a_{3,3}  + {\rm signe} \left( \begin{array}{ccc} 1 & 2 & 3 \\2 & 3 & 1 \end{array} \right) \cdot a_{1,2} a_{2,3} a_{3,1}  \\
        &   & \quad + {\rm signe} \left( \begin{array}{ccc} 1 & 2 & 3 \\3 & 1 & 2 \end{array} \right) \cdot a_{1,3} a_{2,1} a_{3,2}  + {\rm signe} \left( \begin{array}{ccc} 1 & 2 & 3 \\3 & 2 & 1 \end{array} \right) \cdot a_{1,3} a_{2,2} a_{3,1}  \\\\
        & = & (+1) \cdot a_{1,1} a_{2,2} a_{3,3}  + (-1) \cdot a_{1,1} a_{2,3} a_{3,2}  \\
        &   & \quad + (-1) \cdot a_{1,2} a_{2,1} a_{3,3}  + (+1) \cdot a_{1,2} a_{2,3} a_{3,1}  \\
        &   & \quad + (+1) \cdot a_{1,3} a_{2,1} a_{3,2}  + (-1) \cdot a_{1,3} a_{2,2} a_{3,1} 
        \\\\
        & = & 1 \cdot 1 \cdot 2 - 1 \cdot 2 \cdot 3 - 3 \cdot 3 \cdot 2 + 3 \cdot 2 \cdot 2 + 2 \cdot 3 \cdot 3 - 2 \cdot 1 \cdot 2
        \\\\
        & = & 2 - 6 - 18 + 12 + 18 - 4 \\\\
        & = & 4
\end{array}\end{displaymath}%
\lthtmldisplayZ
\lthtmlcheckvsize\clearpage}

{\newpage\clearpage
\lthtmlinlinemathA{tex2html_wrap_inline1278}%
$ {\rm cof} A = \left( \begin{array}{rrr} -4 & -2 & 7 \\0 &
-2 & 3 \\4 & 4 & -8 \end{array} \right)$%
\lthtmlinlinemathZ
\lthtmlcheckvsize\clearpage}

{\newpage\clearpage
\lthtmlinlinemathA{tex2html_wrap_indisplay1280}%
$\displaystyle A^{-1} = \frac{1}{4} \left( \begin{array}{rrr} -4 & 0 & 4
\\-2 & -2 & 4 \\7 & 3 & -8 \end{array} \right).$%
\lthtmlindisplaymathZ
\lthtmlcheckvsize\clearpage}

{\newpage\clearpage
\lthtmldisplayA{displaymath1282}%
\begin{displaymath}
\det A =
\left|
\begin{array}{ccc}
1 & 3 & 2 \\
3 & 1 & 2 \\
2 & 3 & 2
\end{array}
\right|
= 2 + 12 + 18 - 4 - 6 - 18 = 4
\end{displaymath}%
\lthtmldisplayZ
\lthtmlcheckvsize\clearpage}

{\newpage\clearpage
\lthtmldisplayA{displaymath1284}%
\begin{displaymath}
\Delta_1 =
\left|
\begin{array}{ccc}
3 & 3 & 2 \\
4 & 1 & 2 \\
4 & 3 & 2
\end{array}
\right|
= 6 + 24 + 24 - 8 - 18 - 24 = 4
\end{displaymath}%
\lthtmldisplayZ
\lthtmlcheckvsize\clearpage}

{\newpage\clearpage
\lthtmldisplayA{displaymath1286}%
\begin{displaymath}
\Delta_2 =
\left|
\begin{array}{ccc}
1 & 3 & 2 \\
3 & 4 & 2 \\
2 & 4 & 2
\end{array}
\right|
= 8 + 12 + 24 - 16 - 8 - 18 = 2
\end{displaymath}%
\lthtmldisplayZ
\lthtmlcheckvsize\clearpage}

{\newpage\clearpage
\lthtmldisplayA{displaymath1288}%
\begin{displaymath}
\Delta_3 =
\left|
\begin{array}{ccc}
1 & 3 & 3 \\
3 & 1 & 4 \\
2 & 3 & 4
\end{array}
\right|
= 4 + 24 + 27 - 6 - 12 - 36 = 1
\end{displaymath}%
\lthtmldisplayZ
\lthtmlcheckvsize\clearpage}

{\newpage\clearpage
\lthtmldisplayA{displaymath1290}%
\begin{displaymath}
\left(
\begin{array}{c}
x_1 \\
x_2 \\
x_3
\end{array}
\right)
=
\frac{1}{4}
\left(
\begin{array}{r}
4 \\
2 \\
1
\end{array}
\right)
=
\left(
\begin{array}{r}
1 \\
\frac{1}{2} \\
\frac{1}{4}
\end{array}
\right)
\end{displaymath}%
\lthtmldisplayZ
\lthtmlcheckvsize\clearpage}

{\newpage\clearpage
\lthtmldisplayA{displaymath1297}%
\begin{displaymath}\begin{array}{rl}
A = &   \left(
        \begin{array}{rrr|r}
        1 & 3 & -2 & 3 \\
        2 & 4 & 0  & 8 \\
        3 & 6 & 5  & 17
        \end{array}
        \right)
\\\\
\overset{{\tiny\textrm{$\begin{array}{c} L_2-2L_1 \\L_3-3L_1 \end{array}$}}}{\sim} &    \left(
                                        \begin{array}{rrr|r}
                                        1 & 3 & -2 & 3 \\
                                        0 & -2 & 4  & 2 \\
                                        0 & -3 & 11  & 8
                                        \end{array}
                                        \right)
\\\\
\overset{{\tiny\textrm{$\begin{array}{c} L_3 - \frac{3}{2}L_2 \end{array}$}}}{\sim} &    \left(
                                        \begin{array}{rrr|r}
                                        1 & 3 & -2 & 3 \\
                                        0 & 1 & -2  & -1 \\
                                        0 & 0 &  5  & 5
                                        \end{array}
                                        \right)
\\\\
\end{array}\end{displaymath}%
\lthtmldisplayZ
\lthtmlcheckvsize\clearpage}

{\newpage\clearpage
\lthtmldisplayA{displaymath1299}%
\begin{displaymath}\begin{array}[t]{rclclcl}
x_3 & = & 1 \\\\
x_2 & = & -1 + 2x_3         & = & -1 + 2 \cdot 1            & = & 1 \\\\
x_1 & = & 3 - 3x_2 + 2x_3   & = & 3 - 3 \cdot 1 + 2 \cdot 1 & = & 2.
\end{array}\end{displaymath}%
\lthtmldisplayZ
\lthtmlcheckvsize\clearpage}

{\newpage\clearpage
\lthtmldisplayA{displaymath1301}%
\begin{displaymath}
\left(
\begin{array}{c}
x_1 \\
x_2 \\
x_3
\end{array}
\right)
=
\left(
\begin{array}{c}
2 \\
1 \\
1
\end{array}
\right)
\end{displaymath}%
\lthtmldisplayZ
\lthtmlcheckvsize\clearpage}

{\newpage\clearpage
\lthtmlinlinemathA{tex2html_wrap_inline1303}%
$ m=4$%
\lthtmlinlinemathZ
\lthtmlcheckvsize\clearpage}

{\newpage\clearpage
\lthtmldisplayA{displaymath1305}%
\begin{displaymath}\begin{array}{rll}
A = &   \left(
        \begin{array}{rrr|r}
        1 & 3 & -2 & 3 \\
        2 & 4 & 0  & 8 \\
        3 & 6 & 5  & 17
        \end{array}
        \right)
    & {\rm Ligne} = (3,\cdot,\cdot)
\\\\
\overset{{\tiny\textrm{$\begin{array}{c} L_1-0,333L_3 \\L_2-0,667L_3 \end{array}$}}}{\sim} &    \left(
                                                \begin{array}{rrr|r}
                                                0,000 & 1,000 & -3,666 & -2,666 \\
                                                0,000 & 0,000 & -3,334 & -3,330 \\
                                                3,000 & 6,000 & 5,000  & 17,00
                                                \end{array}
                                                \right)
    & {\rm Ligne} = (3, 1, \cdot)
\\\\
\overset{{\tiny\textrm{$\begin{array}{c}  \end{array}$}}}{\sim} &    \left(
                    \begin{array}{rrr|r}
                    0,000 & 1,000 & -3,666 & -2,666 \\
                    0,000 & 0,000 & -3,334 & -3,330 \\
                    3,000 & 6,000 & 5,000  & 17,00
                    \end{array}
                    \right)
    & {\rm Ligne} = (3, 1, 2)
\\\\
\end{array}\end{displaymath}%
\lthtmldisplayZ
\lthtmlcheckvsize\clearpage}

{\newpage\clearpage
\lthtmldisplayA{displaymath1307}%
\begin{displaymath}\begin{array}[t]{rclclcl}
x_3 & = & -3,330/-3,334              & = & 0,9988 \\\\
x_2 & = & -2,666 + 3,666x_3          & = & -2,666 + 3,662       & = & 0,996 \\\\
x_1 & = & \frac{17,00 - 6,000x_2 - 5,000x_3}{3,000}  & = & \frac{17,00 - 5,976 - 4,994}{3,00} & = & \frac{6,026}{3} = 2,009.
\end{array}\end{displaymath}%
\lthtmldisplayZ
\lthtmlcheckvsize\clearpage}

{\newpage\clearpage
\lthtmldisplayA{displaymath1309}%
\begin{displaymath}
\left(
\begin{array}{c}
x_1 \\
x_2 \\
x_3
\end{array}
\right)
=
\left(
\begin{array}{c}
2,009 \\
0,996 \\
0,9988
\end{array}
\right)
\end{displaymath}%
\lthtmldisplayZ
\lthtmlcheckvsize\clearpage}

{\newpage\clearpage
\lthtmldisplayA{displaymath1313}%
\begin{displaymath}\begin{array}{rll}
A = &   \left(
        \begin{array}{rrr|r}
        1 & 3 & -2 & 3 \\
        2 & 4 & 0  & 8 \\
        3 & 6 & 5  & 17
        \end{array}
        \right)
    & \begin{array}{l} {\rm Ligne} = (3,\cdot,\cdot) \\{\rm Colonne} = (2,\cdot,\cdot) \end{array}
\\\\
\overset{{\tiny\textrm{$\begin{array}{c} L_1-0,5L_3 \\L_2-0,667L_3 \end{array}$}}}{\sim} &    \left(
                                                \begin{array}{rrr|r}
                                                -0,5000 & 0,000 & -4,500 & -5,500 \\
                                                0,000   & 0,000 & -3,334 & -3,330 \\
                                                3,000   & 6,000 & 5,000  & 17,00
                                                \end{array}
                                                \right)
    & \begin{array}{l} {\rm Ligne} = (3,1,\cdot) \\{\rm Colonne} = (2,3,\cdot) \end{array}
\\\\
\overset{{\tiny\textrm{$\begin{array}{c} L_2-\frac{3,34}{4,50}L_1 \end{array}$}}}{\sim} &    \left(
                    \begin{array}{rrr|r}
                    -0,5000 & 0,000 & -4,500 & -5,500 \\
                    0,3705  & 0,000 & 0,000  & 0,7450 \\
                    3,000   & 6,000 & 5,000  & 17,00
                    \end{array}
                    \right)
    & \begin{array}{l} {\rm Ligne} = (3,1,2) \\{\rm Colonne} = (2,3,1) \end{array}
\\\\
\end{array}\end{displaymath}%
\lthtmldisplayZ
\lthtmlcheckvsize\clearpage}

{\newpage\clearpage
\lthtmldisplayA{displaymath1315}%
\begin{displaymath}\begin{array}[t]{rclclcl}
x_1 & = & 0,7450/0,3705                            & = & 2,011 \\\\
x_3 & = & \frac{-5,500 + 0,5000x_1}{-4,500}         & = & \frac{-5,500 + 1,006}{-4,500}       & = & 0,9987 \\\\
x_2 & = & \frac{17,00 - 3,000x_1 - 5,000x_3}{6,000}  & = & \frac{17,00 - 6,033 - 4,994}{6,000} & = & \frac{5,970}{6,000} = 0,995.
\end{array}\end{displaymath}%
\lthtmldisplayZ
\lthtmlcheckvsize\clearpage}

{\newpage\clearpage
\lthtmldisplayA{displaymath1317}%
\begin{displaymath}
\left(
\begin{array}{c}
x_1 \\
x_2 \\
x_3
\end{array}
\right)
=
\left(
\begin{array}{c}
2,011 \\
0,995 \\
0,9987
\end{array}
\right)
\end{displaymath}%
\lthtmldisplayZ
\lthtmlcheckvsize\clearpage}

{\newpage\clearpage
\lthtmlinlinemathA{tex2html_wrap_indisplay1319}%
$\displaystyle \begin{array}{rcl}
r(x_0)  & = & b-Ax_0 \\
& = &
\left(
\begin{array}{c}
3 \\8 \\17
\end{array}
\right)
-
\left(
\begin{array}{rrr}
1 & 3 & -2 \\
2 & 4 & 0  \\
3 & 6 & 5
\end{array}
\right)
\left(
\begin{array}{c}
2 \\1 \\1
\end{array}
\right) \\
& = &
\left(
\begin{array}{c}
3 \\8 \\17
\end{array}
\right)
-
\left(
\begin{array}{c}
3 \\8 \\17
\end{array}
\right) \\
& = &
\left(
\begin{array}{c}
0 \\0 \\0
\end{array}
\right).
\end{array}
$%
\lthtmlindisplaymathZ
\lthtmlcheckvsize\clearpage}

{\newpage\clearpage
\lthtmlinlinemathA{tex2html_wrap_indisplay1321}%
$\displaystyle \begin{array}{rcl}
r(x_0)  & = & b-Ax_0 \\
& = &
\left(
\begin{array}{c}
3 \\8 \\17
\end{array}
\right)
-
\left(
\begin{array}{rrr}
1 & 3 & -2 \\
2 & 4 & 0  \\
3 & 6 & 5
\end{array}
\right)
\left(
\begin{array}{c}
2,007 \\0,997 \\0,9988
\end{array}
\right) \\
& = &
\left(
\begin{array}{c}
3 \\8 \\17
\end{array}
\right)
-
\left(
\begin{array}{c}
3,000 \\8,002 \\17,00
\end{array}
\right) \\
& = &
\left(
\begin{array}{c}
0,000 \\-0,002000 \\0,000
\end{array}
\right).
\end{array}
$%
\lthtmlindisplaymathZ
\lthtmlcheckvsize\clearpage}

{\newpage\clearpage
\lthtmlinlinemathA{tex2html_wrap_indisplay1323}%
$\displaystyle \begin{array}{rcl}
r(x_0)  & = & b-Ax_0 \\
& = &
\left(
\begin{array}{c}
3 \\8 \\17
\end{array}
\right)
-
\left(
\begin{array}{rrr}
1 & 3 & -2 \\
2 & 4 & 0  \\
3 & 6 & 5
\end{array}
\right)
\left(
\begin{array}{c}
2,011 \\0,995 \\0,9987
\end{array}
\right) \\
& = &
\left(
\begin{array}{c}
3 \\8 \\17
\end{array}
\right)
-
\left(
\begin{array}{c}
2,999 \\8,002 \\16,99
\end{array}
\right) \\
& = &
\left(
\begin{array}{c}
0,001000 \\-0,002000 \\0,01000
\end{array}
\right).
\end{array}
$%
\lthtmlindisplaymathZ
\lthtmlcheckvsize\clearpage}

{\newpage\clearpage
\lthtmldisplayA{displaymath1325}%
\begin{displaymath}\begin{array}{rcl}
    \| r(x_0) \|    & = & \sqrt{(-0,000)^2 + (-0,002000)^2 + (0,000)^2} \\
                    & = & \sqrt{0,000 + 0,000004000 + 0,000} \\
                    & = & \sqrt{0,000004000} \\
                    & = & 0,002000
    \end{array}\end{displaymath}%
\lthtmldisplayZ
\lthtmlcheckvsize\clearpage}

{\newpage\clearpage
\lthtmldisplayA{displaymath1327}%
\begin{displaymath}\begin{array}{rcl}
    \| r(x_0) \|    & = & \sqrt{(0,001000)^2 + (-0,002000)^2 + (0,01000)^2} \\
                    & = & \sqrt{0,0000010000 + 0,000004000 + 0,0001000} \\
                    & = & \sqrt{0,000005000 + 0,0001000} \\
                    & = & \sqrt{0,0001050} \\
                    & = & 0,01025
    \end{array}\end{displaymath}%
\lthtmldisplayZ
\lthtmlcheckvsize\clearpage}

{\newpage\clearpage
\lthtmlinlinemathA{tex2html_wrap_inline1329}%
$ \frac{2}{3}n^3 + \frac{3}{2}n^2 - \frac{7}{6}n$%
\lthtmlinlinemathZ
\lthtmlcheckvsize\clearpage}

{\newpage\clearpage
\lthtmlinlinemathA{tex2html_wrap_inline1331}%
$ A^{(0)}$%
\lthtmlinlinemathZ
\lthtmlcheckvsize\clearpage}

{\newpage\clearpage
\lthtmlinlinemathA{tex2html_wrap_inline1333}%
$ n^2$%
\lthtmlinlinemathZ
\lthtmlcheckvsize\clearpage}

{\newpage\clearpage
\lthtmlinlinemathA{tex2html_wrap_inline1335}%
$ n^2-1$%
\lthtmlinlinemathZ
\lthtmlcheckvsize\clearpage}

{\newpage\clearpage
\lthtmlinlinemathA{tex2html_wrap_inline1337}%
$ A^{(1)}$%
\lthtmlinlinemathZ
\lthtmlcheckvsize\clearpage}

{\newpage\clearpage
\lthtmlinlinemathA{tex2html_wrap_inline1339}%
$ (n-1)^2$%
\lthtmlinlinemathZ
\lthtmlcheckvsize\clearpage}

{\newpage\clearpage
\lthtmlinlinemathA{tex2html_wrap_inline1341}%
$ (n-1)^2 -1$%
\lthtmlinlinemathZ
\lthtmlcheckvsize\clearpage}

{\newpage\clearpage
\lthtmlinlinemathA{tex2html_wrap_inline1343}%
$ A^{(k)}$%
\lthtmlinlinemathZ
\lthtmlcheckvsize\clearpage}

{\newpage\clearpage
\lthtmlinlinemathA{tex2html_wrap_inline1345}%
$ (n-k)^2$%
\lthtmlinlinemathZ
\lthtmlcheckvsize\clearpage}

{\newpage\clearpage
\lthtmlinlinemathA{tex2html_wrap_inline1347}%
$ (n-k)^2 -1$%
\lthtmlinlinemathZ
\lthtmlcheckvsize\clearpage}

{\newpage\clearpage
\lthtmlinlinemathA{tex2html_wrap_inline1349}%
$ A^{n-1}$%
\lthtmlinlinemathZ
\lthtmlcheckvsize\clearpage}

{\newpage\clearpage
\lthtmlinlinemathA{tex2html_wrap_inline1351}%
$ k$%
\lthtmlinlinemathZ
\lthtmlcheckvsize\clearpage}

{\newpage\clearpage
\lthtmlinlinemathA{tex2html_wrap_inline1353}%
$ n-2$%
\lthtmlinlinemathZ
\lthtmlcheckvsize\clearpage}

{\newpage\clearpage
\lthtmldisplayA{displaymath1355}%
\begin{displaymath}\begin{array}{rcl}
\sum\limits_{i=0}^{n-2} [ (n-k)^2 -1 ]
    & = & \sum\limits_{k=0}^{n-2} (n-k)^2 - \sum\limits_{i=0}^{n-2} 1 \\\\
    & = & \sum\limits_{k=0}^{n-2} (n-k)^2 - (n-1) \\\\
    & = & \sum\limits_{k=2}^{n} k^2 - (n-1) \\\\
    & = & \sum\limits_{k=1}^{n} k^2 - 1 - (n-1) \\\\
    & = & \sum\limits_{k=1}^{n} k^2 - n \\\\
    & = & \frac{n(n+1)(2n+1)}{6} - n \\\\
    & = & \frac{2n^3 + 3n^2 + n}{6} - n \\\\
    & = & \frac{2n^3 + 3n^2 -5n}{6} \\\\
    & = & \frac{1}{3}n^3 + \frac{1}{2}n^2 -\frac{5}{6}n
\end{array}\end{displaymath}%
\lthtmldisplayZ
\lthtmlcheckvsize\clearpage}

{\newpage\clearpage
\lthtmlinlinemathA{tex2html_wrap_inline1357}%
$ \frac{2}{3}n^3 + \frac{3}{2}n^2 - \frac{7}{6}n +
\frac{1}{3}n^3 + \frac{1}{2}n^2 -\frac{5}{6}n =
n^3 + 2n^2 -2n$%
\lthtmlinlinemathZ
\lthtmlcheckvsize\clearpage}

{\newpage\clearpage
\lthtmlinlinemathA{tex2html_wrap_inline1359}%
$ t=0$%
\lthtmlinlinemathZ
\lthtmlcheckvsize\clearpage}

{\newpage\clearpage
\lthtmldisplayA{displaymath1361}%
\begin{displaymath}
\begin{array}{rcrcrcrcr}
x(t) & = & 4Ae^{-t} & + & 14Be^{-2t} & + & Ce^{-3t} & + & e^{t} \\
y(t) & = & -9Ae^{-t} & - & 30Be^{-2t} & - & 2Ce^{-3t} &  &  \\
z(t) & = & 6Ae^{-t} & + & 17Be^{-2t} & + & Ce^{-3t} & - & e^{t}
\end{array}
\end{displaymath}%
\lthtmldisplayZ
\lthtmlcheckvsize\clearpage}

{\newpage\clearpage
\lthtmldisplayA{displaymath1363}%
\begin{displaymath}
\begin{array}{rcrcrcrcr}
0 & = & 4A  & + & 14B & + & C  & + & 1  \\
0 & = & -9A & - & 30B & - & 2C &   &    \\
0 & = & 6A  & + & 17B & + & C  & - & 1
\end{array}
\end{displaymath}%
\lthtmldisplayZ
\lthtmlcheckvsize\clearpage}

{\newpage\clearpage
\lthtmldisplayA{displaymath1365}%
\begin{displaymath}\begin{array}{rll}
A = &   \left(
        \begin{array}{rrr|r}
        4  & 14  & 1  & -1 \\
        -9 & -30 & -2 & 0  \\
        6  & 17  & 1  & 1
        \end{array}
        \right)
    & \begin{array}{l} {\rm Ligne} = (2,\cdot,\cdot) \\{\rm Colonne} = (2,\cdot,\cdot) \end{array}
\\\\
\overset{{\tiny\textrm{$\begin{array}{c} L_1+\frac{14}{30}L_2 \\L_3+\frac{17}{30}L_2 \end{array}$}}}{\sim}
    &    \left(
         \begin{array}{rrr|r}
         \frac{-6}{30}  &  0  & \frac{2}{30}    & -1 \\
        -9              & -30 & -2              & 0  \\
         \frac{27}{30}   &  0  & \frac{-4}{30}   & 1
\par
\end{array}
         \right)
    & \begin{array}{l} {\rm Ligne} = (2,3,\cdot) \\{\rm Colonne} = (2,1,\cdot) \end{array}
\\\\
\overset{{\tiny\textrm{$\begin{array}{c} L_1+\frac{2}{9}L_3 \end{array}$}}}{\sim} &
        \left(
        \begin{array}{rrr|r}
         0              &  0  & \frac{1}{27}   & \frac{-7}{9}  \\
        -9              & -30 & -2              & 0             \\
        \frac{9}{10}    &  0  & \frac{-2}{15}   & 1
        \end{array}
        \right)
    & \begin{array}{l} {\rm Ligne} = (2,3,1) \\{\rm Colonne} = (2,1,3) \end{array}
\\\\
\end{array}\end{displaymath}%
\lthtmldisplayZ
\lthtmlcheckvsize\clearpage}

{\newpage\clearpage
\lthtmldisplayA{displaymath1367}%
\begin{displaymath}\begin{array}[t]{rclclcl}
x_3 & = & -21 \\\\
x_1 & = & \frac{1 + \frac{2}{15}x_3}{\frac{9}{10}}  & = & \frac{\frac{-27}{15}}{\frac{9}{10}}    & = & -2 \\\\
x_2 & = & \frac{9x_1 +2x_3}{-30}                    & = & \frac{-60}{-30}           & = & 2.
\end{array}\end{displaymath}%
\lthtmldisplayZ
\lthtmlcheckvsize\clearpage}

{\newpage\clearpage
\lthtmldisplayA{displaymath1369}%
\begin{displaymath}
\begin{array}{rcrcrcrcr}
x(t) & = & -8e^{-t}  & + & 28e^{-2t} & + & -21e^{-3t} & + & e^{t} \\
y(t) & = & 18e^{-t} & - & 60e^{-2t} & + & 42e^{-3t} &   &       \\
z(t) & = & -12e^{-t}  & + & 34e^{-2t}  & + & -21e^{-3t} & - & e^{t}.
\end{array}
\end{displaymath}%
\lthtmldisplayZ
\lthtmlcheckvsize\clearpage}

{\newpage\clearpage
\lthtmlinlinemathA{tex2html_wrap_inline1374}%
$ (T|b)$%
\lthtmlinlinemathZ
\lthtmlcheckvsize\clearpage}

{\newpage\clearpage
\lthtmlinlinemathA{tex2html_wrap_inline1376}%
$ T^{(0)}=T$%
\lthtmlinlinemathZ
\lthtmlcheckvsize\clearpage}

{\newpage\clearpage
\lthtmlinlinemathA{tex2html_wrap_inline1378}%
$ b^{(0)}=b$%
\lthtmlinlinemathZ
\lthtmlcheckvsize\clearpage}

{\newpage\clearpage
\lthtmlinlinemathA{tex2html_wrap_inline1380}%
$ k=1,2,\ldots,n-1$%
\lthtmlinlinemathZ
\lthtmlcheckvsize\clearpage}

{\newpage\clearpage
\lthtmlinlinemathA{tex2html_wrap_inline1382}%
$ T_{k,k}^{(k-1)} = 0$%
\lthtmlinlinemathZ
\lthtmlcheckvsize\clearpage}

{\newpage\clearpage
\lthtmlinlinemathA{tex2html_wrap_inline1386}%
$ k+1$%
\lthtmlinlinemathZ
\lthtmlcheckvsize\clearpage}

{\newpage\clearpage
\lthtmlinlinemathA{tex2html_wrap_inline1388}%
$ T^{(k-1)}$%
\lthtmlinlinemathZ
\lthtmlcheckvsize\clearpage}

{\newpage\clearpage
\lthtmlinlinemathA{tex2html_wrap_inline1390}%
$ b^{(k-1)}$%
\lthtmlinlinemathZ
\lthtmlcheckvsize\clearpage}

{\newpage\clearpage
\lthtmlinlinemathA{tex2html_wrap_inline1392}%
$ m_{k+1,k} = \frac{T_{k+1,k}^{(k-1)}}{T_{k,k}^{(k-1)}}$%
\lthtmlinlinemathZ
\lthtmlcheckvsize\clearpage}

{\newpage\clearpage
\lthtmlinlinemathA{tex2html_wrap_inline1394}%
$ T_{k+1,k+1}^{(k)} = T_{k+1,k+1}^{(k-1)} - m_{k+1,k}T_{k,k+1}^{(k-1)}$%
\lthtmlinlinemathZ
\lthtmlcheckvsize\clearpage}

{\newpage\clearpage
\lthtmlinlinemathA{tex2html_wrap_inline1396}%
$ b_{k+1}^{(k)} = b_{k+1}^{(k-1)} - m_{k+1,k}b_{k}^{(k-1)}$%
\lthtmlinlinemathZ
\lthtmlcheckvsize\clearpage}

{\newpage\clearpage
\lthtmlinlinemathA{tex2html_wrap_inline1415}%
$ T'x=b'$%
\lthtmlinlinemathZ
\lthtmlcheckvsize\clearpage}

{\newpage\clearpage
\lthtmlinlinemathA{tex2html_wrap_inline1417}%
$ b'= (b_1',b_2',\ldots,b_n')^{t}$%
\lthtmlinlinemathZ
\lthtmlcheckvsize\clearpage}

{\newpage\clearpage
\lthtmldisplayA{displaymath1419}%
\begin{displaymath}
T' =
\left(
\begin{array}{ccccccccccc}
d_1'     & u_1       & 0         & 0         & 0         & \cdots & 0        & 0         & 0         & 0       & 0         \\
0       & d_2'      & u_2       & 0         & 0         & \cdots & 0        & 0         & 0         & 0       & 0         \\
0       & 0         & d_3'      & u_3       & 0         & \cdots & 0        & 0         & 0         & 0       & 0         \\
0       & 0         & 0         & d_4'      & u_4       & \cdots & 0        & 0         & 0         & 0       & 0         \\
0       & 0         & 0         & 0         & d_5'      & \ddots & 0        & 0         & 0         & 0       & 0         \\
0       & 0         & 0         & 0         & 0         & \ddots & \ddots   & 0         & 0         & 0       & 0         \\
\vdots  & \vdots    & \vdots    & \vdots    & \vdots    & \ddots & \ddots   & \ddots    & \vdots    & \vdots  & \vdots    \\
0       & 0         & 0         & 0         & 0         & \cdots & 0        & d_{n-3}'  & u_{n-3}   & 0       & 0         \\
0       & 0         & 0         & 0         & 0         & \cdots & 0        & 0         & d_{n-2}'  & u_{n-2} & 0         \\
0       & 0         & 0         & 0         & 0         & \cdots & 0        & 0         & 0         & d_{n-1}' & u_{n-1}   \\
0       & 0         & 0         & 0         & 0         & \cdots & 0        & 0         & 0         & 0       & d_{n}'
\end{array}
\right)
\end{displaymath}%
\lthtmldisplayZ
\lthtmlcheckvsize\clearpage}

{\newpage\clearpage
\lthtmlinlinemathA{tex2html_wrap_inline1421}%
$ x_n = \frac{b_n'}{d_n'}$%
\lthtmlinlinemathZ
\lthtmlcheckvsize\clearpage}

{\newpage\clearpage
\lthtmlinlinemathA{tex2html_wrap_inline1423}%
$ x_i = \frac{b_i' - u_{i}x_{i+1}}{d_i'}$%
\lthtmlinlinemathZ
\lthtmlcheckvsize\clearpage}

{\newpage\clearpage
\lthtmlinlinemathA{tex2html_wrap_inline1425}%
$ i = n-1, n-2, \ldots, 1$%
\lthtmlinlinemathZ
\lthtmlcheckvsize\clearpage}

{\newpage\clearpage
\lthtmlinlinemathA{tex2html_wrap_inline1427}%
$ m_{k+1,k}$%
\lthtmlinlinemathZ
\lthtmlcheckvsize\clearpage}

{\newpage\clearpage
\lthtmlinlinemathA{tex2html_wrap_inline1429}%
$ T_{k+1,k+1}^{(k)}$%
\lthtmlinlinemathZ
\lthtmlcheckvsize\clearpage}

{\newpage\clearpage
\lthtmlinlinemathA{tex2html_wrap_inline1431}%
$ b_{k+1}^{(k)}$%
\lthtmlinlinemathZ
\lthtmlcheckvsize\clearpage}

{\newpage\clearpage
\lthtmlinlinemathA{tex2html_wrap_inline1435}%
$ 5(n-1)$%
\lthtmlinlinemathZ
\lthtmlcheckvsize\clearpage}

{\newpage\clearpage
\lthtmlinlinemathA{tex2html_wrap_inline1437}%
$ x_n$%
\lthtmlinlinemathZ
\lthtmlcheckvsize\clearpage}

{\newpage\clearpage
\lthtmlinlinemathA{tex2html_wrap_inline1439}%
$ x_{n-1}, \ldots,x_2,x_1$%
\lthtmlinlinemathZ
\lthtmlcheckvsize\clearpage}

{\newpage\clearpage
\lthtmlinlinemathA{tex2html_wrap_inline1441}%
$ 3(n-1) +1$%
\lthtmlinlinemathZ
\lthtmlcheckvsize\clearpage}

{\newpage\clearpage
\lthtmlinlinemathA{tex2html_wrap_inline1443}%
$ 5(n-1) + 3(n-1) + 1 = 8n-7$%
\lthtmlinlinemathZ
\lthtmlcheckvsize\clearpage}

{\newpage\clearpage
\lthtmldisplayA{displaymath1445}%
\begin{displaymath}\begin{array}{rl}
A = &   \left(
        \begin{array}{rrrr|r}
        5  & -4 & 0  & 0  & 0 \\
        -4 & 6  & -4 & 0  & 0 \\
        0  & -4 & 6  & -4 & 1 \\
        0  & 0  & -4 & 5  & 0
        \end{array}
        \right)
\\\\
\overset{{\tiny\textrm{$\begin{array}{c} L_2+\frac{4}{5}L_1 \end{array}$}}}{\sim}
    &   \left(
        \begin{array}{rrrr|r}
        5  & -4             & 0  & 0  & 0 \\
        0  & \frac{14}{5}   & -4 & 0  & 0 \\
        0  & -4             & 6  & -4 & 1 \\
        0  & 0              & -4 & 5  & 0
        \end{array}
        \right)
\\\\
\overset{{\tiny\textrm{$\begin{array}{c} L_3 + \frac{10}{7}L_2 \end{array}$}}}{\sim}
    &   \left(
        \begin{array}{rrrr|r}
        5  & -4             & 0             & 0  & 0 \\
        0  & \frac{14}{5}   & -4            & 0  & 0 \\
        0  & 0              & \frac{2}{7}   & -4 & 1 \\
        0  & 0              & -4            & 5  & 0
        \end{array}
        \right)
\\\\
\overset{{\tiny\textrm{$\begin{array}{c} L_4 + 14L_3 \end{array}$}}}{\sim}
    &   \left(
        \begin{array}{rrrr|r}
        5  & -4             & 0             & 0   & 0 \\
        0  & \frac{14}{5}   & -4            & 0   & 0 \\
        0  & 0              & \frac{2}{7}   & -4  & 1 \\
        0  & 0              & 0             & -51 & 14
        \end{array}
        \right)
\\\\
\end{array}\end{displaymath}%
\lthtmldisplayZ
\lthtmlcheckvsize\clearpage}

{\newpage\clearpage
\lthtmldisplayA{displaymath1447}%
\begin{displaymath}\begin{array}[t]{rclclcl}
x_4 & = & \frac{-14}{51} \\\\
x_3 & = & \frac{1 +4 x_4}{\frac{2}{7}}      & = & \frac{\frac{-5}{51}}{\frac{2}{7}}        & = & \frac{-35}{102} \\\\
x_2 & = & \frac{4x_3}{\frac{14}{5}}         & = & \frac{\frac{-70}{51}}{\frac{14}{5}}      & = & \frac{-25}{51} \\\\
x_1 & = & \frac{4x_2}{5}                    & = & \frac{\frac{-100}{51}}{5}                & = & \frac{-20}{51}.
\end{array}\end{displaymath}%
\lthtmldisplayZ
\lthtmlcheckvsize\clearpage}

{\newpage\clearpage
\lthtmldisplayA{displaymath1449}%
\begin{displaymath}
\left(
\begin{array}{c}
x_1 \\
x_2 \\
x_3 \\
x_4
\end{array}
\right)
=
\left(
\begin{array}{c}
\frac{-20}{51} \\
\frac{-25}{51} \\
\frac{-35}{102} \\
\frac{-14}{51}
\end{array}
\right)
\end{displaymath}%
\lthtmldisplayZ
\lthtmlcheckvsize\clearpage}


\end{document}
